% Nouveau EGA V (cote 25): archive cote 25, site page 63 (PDF page 64).
% typescript. Draft compared: none.
% AI-transcribed by Claude Opus 5.5 (claude-opus-5-5, adaptive thinking, effort high), Message Batches API, batch ega25b (msgbatch_01A7n5YuXubUPNK1zsepi8RG), 2026-11-13; not yet checked by eye.
\scanpage{25/63}
\mspage{58}

apparait comme un sous-préschéma fermé de $P(\overset{\hand{\unsure{n}}}{\bigwedge}\underline{E})$; en particulier, il
est séparé sur S et quasi-compact sur S, et  si $\underline{E}$  est de type fini,
il est projectif sur S. Si $\underline{E}$ est de \unsure{présentation|présantation} finie, alors
il en est de même de $\underline{\text{Grass}}_n(\underline{E})$ : en effet, \xout{\unsure{localement sur S,}} on peut
supposer S affine, S=Spec(A), alors $\underline{E}$ provient d'un module de type fini
sur un sous-anneau de type fini de A, -  or la formation de $\underline{\text{Grass}}_n(\underline{E})$
est évidemment compatible avec la changement de base sur S.

Lorsque $\underline{E}$  est localement libre,\xout{de type fini,} alors  $\underline{\text{Grass}}_n(\underline{E})$
est \uline{lisse} sur S\inserted{\unsure{(à fibres géom. connexes}}. Cela provient du fait plus précis: si $\underline{E}$ est libre,
de rang r, alors  $\underline{\text{Grass}}_n(\underline{E})$  peut être recouvert par  $\binom{r}{n}$  ouverts
\inserted{S-}isomorphes chacun  à  l'espace affine type de dimension relative
n(r-n)  sur S.  Cette décomposition correspond au choix grâce à la base
de $\underline{E}$  de  $\binom{r}{n}$  décompositions de $\underline{E}$  en \xout{sommes directes} suites
exactes
\[
\text{(s)} \qquad 0 \rightarrow \underline{E}' \relbar \underline{E} \relbar \underline{E}'' \longrightarrow 0 \quad ,
\]
avec $\underline{E}'$  loc libre de rang n. \xout{\illegible{1}} Une telle suite exacte permet
de définir un sous-foncteur  $\underline{\text{Grass}}_n(\text{s})$  de  $\underline{\text{Grass}}_n(\underline{E})$ , en se bornant
aux  quotients  $\underline{G}$  loc libres de rang  n de  $\underline{E}_{S'}$  tels que l'homomorphi-%
sme composé  $\underline{E}'_{S'} \longrightarrow \underline{E}_{S'} \longrightarrow \underline{G}$  soit surjectif (donc bijectif). Or
l'inclusion  $\underline{\text{Grass}}_n(\text{s}) \rightarrow \underline{\text{Grass}}_n(\underline{E})$  est représentable \uline{par immersions
ouvertes}, et d'autre\xout{s} part  $\underline{\text{Grass}}_n(\text{s})$  est canoniquement isomorphe
au fibré ve\unsure{c}toriel  \xout{\unsure{V(E''⊗E')}}  V ( $\underline{\text{Hom}}_{\underline{O}_S}(\underline{E}',\underline{E}'')$) .

Comme conséquence par exemple de cette structure particulière,
on peut signaler que si s $\in$ S, alors ($\underline{E}$ étant loc libre de rang fini)
tout point de $\underline{\text{Grass}}_n(\underline{E})$  à valeurs dans k(s)  se relève en une section
au dessus d'un voisinage de s.  D'autre part, si S=Spec(k), \xout{\illegible{1}} k
un corps fnfini, alors tout ouvert non vide de $\underline{\text{Grassn}}(\underline{E})$  contient un

\mspage{63}
